% =============================================================================
\bloco{1}{Exam level: the three approaches, item by item}{Kurlat \S1.1; Jones ch.~2 \,\textbullet\, 30 points}

\questao{(20 points) \textit{Source: Kurlat, P. (2020), \emph{A Course in Modern Macroeconomics}, Exercise 1.1 ``Accounting'',
p.~27 --- statement reproduced verbatim, all items (a)--(g).}}
\textit{What this trains:} building the Production\,$|$\,Income\,$|$\,Expenditure table for imports, inventories,
government output valued at cost, a retail margin, investment goods and taxes --- the exact skill Prova~01
Q1(b)--(c) asked for.

\noindent How does GDP accounting record the following events? For each of them, describe how they would be computed in GDP
accounts using the income method, the production method and the expenditure method.
\begin{enumerate}[label=(\alph*),leftmargin=2.2em,topsep=2pt,itemsep=1pt]
\item A car manufacturer buys components from Japan for \$1 to be used in production later on and stores them at its warehouse.
\item A car manufacturer buys components from Japan for \$1 and uses half of those components in the production of a car, that it
sells to Andy for \$2. It stores the rest of the components.
\item An army battalion is deployed to the border to repel a threatened Canadian invasion. The soldiers earn wages of \$10,000 and
use ammunition that the government buys for \$5,000. The ammunition is produced using \$2,000 of imported steel and 100 hours of
work, for which the workers were paid \$1,000.
\item Walmart sells 1000 bottles of Coca-Cola for \$1,500. It had previously paid \$1,200 for them.
\item A shipyard builds a cruise ship. It pays wages of \$200,000, interest on loans (from US residents) of \$100,000 and \$300,000
for imported raw materials. The ship is sold for \$1,000,000 to a cruise company. In the same year, the cruise company has revenue
for \$50,000 from operating cruises, pays wages of \$20,000 to its workers and has no other expenses. Half the cruise revenue comes
from tourists who reside in the United States and half comes from tourists who reside abroad.
\item The government collects \$1000 in income taxes from Roger.
\item Roger earns \$4000 for working as a babysitter and pays \$1000 in income taxes.
\end{enumerate}
\resolespaco
\begin{sol}
\begin{notacao}
\item $j$: a producer (a firm, a government unit or a self-employed person). $Q_j\ge0$: its gross output, sales plus the goods it
adds to its inventories (minus the goods it draws from them). $Z_j\ge0$: the intermediate inputs it buys in the period.
\item $VA_j$: value added of $j$. $W_j\ge0$, $R_j\ge0$: the wages and the interest it pays. $\Pi_j$: its profit (proprietor's
income for the self-employed), gross of depreciation; it can be zero or negative.
\item $C$, $G$, $X$, $M\ge0$: consumption, government purchases, exports, imports. $I$: investment, including the change in
inventories, so it can be negative.
\item (G1)--(G4): the accounting identities stated in the first step; the later questions of this list cite them.
\end{notacao}

\passo{method}{For every event, list (i) each producer's value added; (ii) the incomes generated, gross of depreciation;
(iii) the final uses. The three columns rest on four identities:}
\begin{align}
\text{GDP} &= \textstyle\sum_j VA_j, \tag{G1}\label{q1:G1}\\
VA_j &= Q_j-Z_j, \tag{G2}\label{q1:G2}\\
\text{GDP} &= \textstyle\sum_j\big(W_j+R_j+\Pi_j\big),\qquad \Pi_j\equiv VA_j-W_j-R_j, \tag{G3}\label{q1:G3}\\
\text{GDP} &= C+I+G+X-M. \tag{G4}\label{q1:G4}
\end{align}
Transfers and taxes enter none of the three.

\passo{(a) components imported and stored}{Inventories rise by 1, so the firm's output is $Q=1$; the imported input is $Z=1$.
Using \eqref{q1:G2}:}
\begin{align}
VA &= (1)-(1) \nonumber\\
&= 0. \label{q1:aVA}
\end{align}
With no wages or interest, \eqref{q1:G3} makes profit equal to \eqref{q1:aVA}:
\begin{align*}
\Pi &= (0)-(0)-(0)\\
&= 0 .
\end{align*}
The final uses are inventory investment $I=1$ and imports $M=1$; by \eqref{q1:G4}:
\begin{align*}
\text{GDP} &= (1)-(1)\\
&= 0 .
\end{align*}
\begin{center}\footnotesize\begin{tabular}{@{}lr@{\qquad}lr@{\qquad}lr@{}}\toprule
\multicolumn{2}{@{}c}{Production} & \multicolumn{2}{c}{Income} & \multicolumn{2}{c@{}}{Expenditure}\\\midrule
Car manufacturing $1-1$ & 0 & Corporate profits & 0 & Investment (inventories) & 1\\
 & & & & Imports & $-1$\\\midrule
Total & 0 & Total & 0 & Total & 0\\\bottomrule\end{tabular}\end{center}
Nothing was produced in the country: the inventory investment is exactly offset by the import.

\passo{(b) half used, half stored}{Output is the car sold plus the components added to inventory:}
\begin{align}
Q &= (2)+(0.5) \nonumber\\
&= 2.5. \label{q1:bQ}
\end{align}
The intermediate input is the imported 1. Using \eqref{q1:G2} with \eqref{q1:bQ}:
\begin{align}
VA &= (2.5)-(1) \nonumber\\
&= 1.5. \label{q1:bVA}
\end{align}
With no wages or interest, \eqref{q1:G3} gives $\Pi=VA$. Directly, profit is revenue minus the cost of the components used up
(the other 0.5 is still an asset, in inventory):
\begin{align*}
\Pi &= (2)-(0.5)\\
&= 1.5 ,
\end{align*}
which matches \eqref{q1:bVA}. The final uses are $C=2$, inventory investment $I=0.5$ and $M=1$; by \eqref{q1:G4}:
\begin{align*}
\text{GDP} &= (2)+(0.5)-(1)\\
&= 1.5 .
\end{align*}
\begin{center}\footnotesize\begin{tabular}{@{}lr@{\qquad}lr@{\qquad}lr@{}}\toprule
\multicolumn{2}{@{}c}{Production} & \multicolumn{2}{c}{Income} & \multicolumn{2}{c@{}}{Expenditure}\\\midrule
Car manufacturing $2+0.5-1$ & 1.5 & Corporate profits & 1.5 & Consumption & 2\\
 & & & & Investment (inventories) & 0.5\\
 & & & & Imports & $-1$\\\midrule
Total & 1.5 & Total & 1.5 & Total & 1.5\\\bottomrule\end{tabular}\end{center}
Confirmed against Kurlat's own worked Example~1.6 (p.~19), the same event at ten times the scale: 15 by all three approaches.

\passo{(c) army battalion}{Government output has no price, so it is valued at cost, wages plus ammunition:}
\begin{align}
Q_{\mathrm{def}} &= (10{,}000)+(5{,}000) \nonumber\\
&= 15{,}000. \label{q1:cQ}
\end{align}
The ammunition is an intermediate input of defence. Using \eqref{q1:G2} with \eqref{q1:cQ}:
\begin{align}
VA_{\mathrm{def}} &= (15{,}000)-(5{,}000) \nonumber\\
&= 10{,}000. \label{q1:cVAd}
\end{align}
The ammunition maker sells 5{,}000 and uses 2{,}000 of imported materials; by \eqref{q1:G2}:
\begin{align}
VA_{\mathrm{am}} &= (5{,}000)-(2{,}000) \nonumber\\
&= 3{,}000. \label{q1:cVAa}
\end{align}
It pays wages of 1{,}000, so by \eqref{q1:G3} its profit is
\begin{align}
\Pi_{\mathrm{am}} &= (5{,}000)-(2{,}000)-(1{,}000) \nonumber\\
&= 2{,}000. \label{q1:cPi}
\end{align}
Product: summing \eqref{q1:cVAd} and \eqref{q1:cVAa} with \eqref{q1:G1},
\begin{align*}
\text{GDP} &= (10{,}000)+(3{,}000)\\
&= 13{,}000 .
\end{align*}
Income: by \eqref{q1:G3}, the soldiers' wages, the ammunition wages and the profit \eqref{q1:cPi},
\begin{align*}
\text{GDP} &= (10{,}000)+(1{,}000)+(2{,}000)\\
&= 13{,}000 .
\end{align*}
Expenditure: $G$ is the cost \eqref{q1:cQ} and $M=2{,}000$; by \eqref{q1:G4},
\begin{align*}
\text{GDP} &= (15{,}000)-(2{,}000)\\
&= 13{,}000 .
\end{align*}
\begin{center}\footnotesize\begin{tabular}{@{}lr@{\qquad}lr@{\qquad}lr@{}}\toprule
\multicolumn{2}{@{}c}{Production} & \multicolumn{2}{c}{Income} & \multicolumn{2}{c@{}}{Expenditure}\\\midrule
Public defence & 10{,}000 & Wages (soldiers) & 10{,}000 & Government & 15{,}000\\
Manufacturing & 3{,}000 & Wages (ammunition) & 1{,}000 & Imports & $-2{,}000$\\
 & & Corporate profits & 2{,}000 & & \\\midrule
Total & 13{,}000 & Total & 13{,}000 & Total & 13{,}000\\\bottomrule\end{tabular}\end{center}

\passo{(d) Walmart}{Reading: the bottles were bought in an earlier period and sat in Walmart's inventory. Walmart sells 1{,}500
and draws 1{,}200 of bottles from inventory, so its output this year is the retail service:}
\begin{align}
Q_W &= (1{,}500)-(1{,}200) \nonumber\\
&= 300. \label{q1:dQ}
\end{align}
It buys no inputs this year, so by \eqref{q1:G2} with \eqref{q1:dQ}, and then by \eqref{q1:G3} with no wages or interest:
\begin{align*}
VA_W &= (300)-(0)\\
&= 300,\\
\Pi_W &= (300)-(0)-(0)\\
&= 300 .
\end{align*}
The final uses are $C=1{,}500$
and inventory investment $I=-1{,}200$ (the bottles leave inventory); by \eqref{q1:G4}:
\begin{align*}
\text{GDP} &= (1{,}500)+(-1{,}200)\\
&= 300 ,
\end{align*}
equal to \eqref{q1:dQ}.
\begin{center}\footnotesize\begin{tabular}{@{}lr@{\qquad}lr@{\qquad}lr@{}}\toprule
\multicolumn{2}{@{}c}{Production} & \multicolumn{2}{c}{Income} & \multicolumn{2}{c@{}}{Expenditure}\\\midrule
Retail trade & 300 & Corporate profits & 300 & Consumption & 1{,}500\\
 & & & & Investment (inventories) & $-1{,}200$\\\midrule
Total & 300 & Total & 300 & Total & 300\\\bottomrule\end{tabular}\end{center}
If instead the bottles were made by a domestic bottler this same year, the bottler's value added (1{,}200) is also this year's, inventories
do not change and, by \eqref{q1:G1} and \eqref{q1:G4}, GDP becomes
\begin{align*}
\text{GDP} &= (1{,}200)+(300)\\
&= 1{,}500=C .
\end{align*}
State the reading you use.

\passo{(e) shipyard and cruise company}{The shipyard sells the ship for 1{,}000{,}000 and uses 300{,}000 of imported materials;
by \eqref{q1:G2}:}
\begin{align}
VA_S &= (1{,}000{,}000)-(300{,}000) \nonumber\\
&= 700{,}000. \label{q1:eVAs}
\end{align}
It pays wages of 200{,}000 and interest of 100{,}000, so by \eqref{q1:G3}:
\begin{align}
\Pi_S &= (1{,}000{,}000)-(200{,}000)-(100{,}000)-(300{,}000) \nonumber\\
&= 400{,}000. \label{q1:ePis}
\end{align}
The cruise company uses no inputs and pays wages of 20{,}000; by \eqref{q1:G2} and \eqref{q1:G3}:
\begin{align}
VA_{Cr} &= (50{,}000)-(0) \nonumber\\
&= 50{,}000, \label{q1:eVAc}\\
\Pi_{Cr} &= (50{,}000)-(20{,}000) \nonumber\\
&= 30{,}000. \label{q1:ePic}
\end{align}
Product: summing \eqref{q1:eVAs} and \eqref{q1:eVAc} with \eqref{q1:G1},
\begin{align*}
\text{GDP} &= (700{,}000)+(50{,}000)\\
&= 750{,}000 .
\end{align*}
Income: by \eqref{q1:G3}, with the profits \eqref{q1:ePis} and \eqref{q1:ePic},
\begin{align*}
\text{GDP} &= \underbrace{\big((200{,}000)+(20{,}000)\big)}_{\text{wages}}+\underbrace{(100{,}000)}_{\text{interest}}
+\underbrace{\big((400{,}000)+(30{,}000)\big)}_{\text{profits}}\\
&= 220{,}000+100{,}000+430{,}000\\
&= 750{,}000 .
\end{align*}
Expenditure: the ship is an investment good, $I=1{,}000{,}000$; the cruises sold to residents are consumption, $C=25{,}000$, and
those sold to non-residents are exports, $X=25{,}000$; $M=300{,}000$. By \eqref{q1:G4}:
\begin{align*}
\text{GDP} &= (1{,}000{,}000)+(25{,}000)+(25{,}000)-(300{,}000)\\
&= 750{,}000 .
\end{align*}
\begin{center}\footnotesize\begin{tabular}{@{}lr@{\qquad}lr@{\qquad}lr@{}}\toprule
\multicolumn{2}{@{}c}{Production} & \multicolumn{2}{c}{Income} & \multicolumn{2}{c@{}}{Expenditure}\\\midrule
Shipbuilding & 700{,}000 & Wages $200{,}000+20{,}000$ & 220{,}000 & Investment & 1{,}000{,}000\\
Cruise services & 50{,}000 & Interest & 100{,}000 & Consumption & 25{,}000\\
 & & Profits $400{,}000+30{,}000$ & 430{,}000 & Exports & 25{,}000\\
 & & & & Imports & $-300{,}000$\\\midrule
Total & 750{,}000 & Total & 750{,}000 & Total & 750{,}000\\\bottomrule\end{tabular}\end{center}

\passo{(f) income tax}{A transfer from Roger to the government: no good or service is produced. Production 0, Income 0, Expenditure 0.}

\passo{(g) babysitting}{Production: household services with no inputs, so by \eqref{q1:G2}; income: proprietor's income
by \eqref{q1:G3} (Roger is self-employed), with no wages or interest:}
\begin{align*}
VA &= (4{,}000)-(0)\\
&= 4{,}000,\\
\Pi &= (4{,}000)-(0)-(0)\\
&= 4{,}000 .
\end{align*}
The 1{,}000 of tax only splits that income between Roger and the government. Expenditure:
consumption 4{,}000 (the parents buy the service), so \eqref{q1:G4} gives the same total. GDP $=4{,}000$.

\intuicao{The three columns are three views of the same production: who made it, who earned it, who used it. When the columns disagree,
something was double-counted (an intermediate input) or left out (inventories, imports, government output at cost).}
\refbib{Kurlat (2020), \S1.1, pp.~15--22: value added (Ex.~1.2, p.~17), inventories (Ex.~1.3, p.~17), imports (Ex.~1.6, p.~19),
government at cost and transfers (Exs.~1.7--1.8, p.~20). No book answer exists for Exercise 1.1; item (b) is checked against Ex.~1.6.}
\rubrica{
\item 3.5 points per item (a)--(e), 1.25 per item (f)--(g); within an item, the three columns weigh equally.
\item Only one approach given, even with the right total: max.\ a third of the item (the Prova~01 Q1(b) error).
\item Imported input or stored components left out of value added: max.\ half of (a)--(b).
\item Taxes counted as income or expenditure in (f)--(g): 0 for those items.
\item Arithmetic slip with the right layout: $-0.5$ point, once per item.
}
\end{sol}

% -----------------------------------------------------------------------------
\questao{(10 points) \textit{Source: Jones, C.\,I. (2020), \emph{Macroeconomics}, 5th ed., ch.~2, Exercise~2 ``What counts as
GDP (I)?'', p.~39, and Exercise~3 ``What counts as GDP (II)?'', pp.~39--40 --- statements reproduced verbatim, all items.}}
\textit{What this trains:} classifying an event in one line --- new production or not, final or intermediate, which
component of $C+I+G+X-M$ moves --- before writing any table.

\noindent\textbf{2. What counts as GDP (I)?} By how much does GDP rise in each of the following scenarios? Explain.
\begin{enumerate}[label=(\alph*),leftmargin=2.2em,topsep=2pt,itemsep=1pt]
\item You spend \$5,000 on college tuition this semester.
\item You buy a used car from a friend for \$2,500.
\item The government spends \$100 million to build a dam.
\item Foreign graduate students work as teaching assistants at the local university and earn \$5,000 each.
\end{enumerate}
\noindent\textbf{3. What counts as GDP (II)?} By how much does GDP rise in each of the following scenarios? Explain.
\begin{enumerate}[label=(\alph*),leftmargin=2.2em,topsep=2pt,itemsep=1pt]
\item A computer company buys parts from a local distributor for \$1 million, assembles the parts, and sells the resulting
computers for \$2 million.
\item A real estate agent sells a house for \$200,000 that the previous owners had bought 10 years earlier for \$100,000. The agent
earns a commission of \$6,000.
\item During a recession, the government raises unemployment benefits by \$100 million.
\item A new U.S. airline purchases and imports \$50 million worth of airplanes from the European company Airbus.
\item A new European airline purchases \$50 million worth of airplanes from the American company Boeing.
\item A store buys \$100,000 of chocolate from Belgium and sells it to consumers in the United States for \$125,000.
\end{enumerate}
\resolespaco
\begin{sol}
\begin{notacao}
\item $\Delta\text{GDP}$: the change in this year's GDP caused by the event; $\Delta C$, $\Delta I$, $\Delta G$, $\Delta X$, $\Delta M$:
the changes in consumption, investment, government purchases, exports and imports (each of any sign).
\item $VA$: value added, as in \eqref{q1:G1}--\eqref{q1:G2} of Question~1.
\end{notacao}
Taking changes in \eqref{q1:G4}, every item below is read off
\begin{equation}
\Delta\text{GDP}=\Delta C+\Delta I+\Delta G+\Delta X-\Delta M. \label{q2:dY}
\end{equation}

\passo{Exercise 2}{(a) $+5{,}000$: education is a service produced this semester, bought by a household ($C$).
(b) $0$: a used car was produced in an earlier year; selling it transfers an existing asset (a dealer's margin, if any, would count).
(c) $+100$ million: new production of a public structure ($G$; government investment).
(d) $+5{,}000$ per student: the work is done inside the country, so it is domestic product whatever the workers' nationality.}

\passo{Exercise 3}{(a) $+2$ million: the final sale. By value added, \eqref{q1:G1} with the parts treated as intermediate:}
\begin{align*}
\Delta\text{GDP} &= \underbrace{(1\text{ million})}_{\text{computer maker}}+\underbrace{(1\text{ million})}_{\text{parts chain}}\\
&= 2\text{ million}.
\end{align*}
(b) $+6{,}000$: only the agent's service is new production; the house and its capital gain are not.
(c) $0$: unemployment benefits are transfers, not purchases of goods or services.
(d) $0$: $I$ rises by 50 million and $M$ by 50 million; by \eqref{q2:dY},
\begin{align*}
\Delta\text{GDP} &= \Delta I-\Delta M\\
&= (50\text{ million})-(50\text{ million})\\
&= 0 .
\end{align*}
(e) $+50$ million: exports, $\Delta X$ in \eqref{q2:dY}.
(f) $+25{,}000$: $C$ rises by 125{,}000 and $M$ by 100{,}000; by \eqref{q2:dY},
\begin{align*}
\Delta\text{GDP} &= \Delta C-\Delta M\\
&= (125{,}000)-(100{,}000)\\
&= 25{,}000 ,
\end{align*}
which is the store's retail margin, the domestic value added.

\intuicao{GDP counts current production of final goods and services inside the country. Each ``0'' above fails one of those words:
not current (used car, house), not production (transfer, capital gain), not domestic (imported aircraft).}
\refbib{Answers confirmed: Exercise~2 against the worked exercise in Jones (2020), p.~41; Exercise~3 against the Jones
\emph{Instructor's Manual} (2016), ch.~2, Exercise~2 (numbered 3 in the 2020 edition), p.~12 --- same six items, same answers.
Kurlat (2020), Ex.~1.4 (financial investment, p.~18) and Ex.~1.8 (transfers, p.~20).}
\rubrica{
\item 1 point per item (ten items), with the component of $C+I+G+X-M$ named; a right number with no reason: half.
\item Used goods, capital gains or transfers counted in GDP: 0 for that item.
\item Imported aircraft counted as a rise in GDP in 3(d): 0 for that item (it is $I$ and $-M$).
}
\end{sol}
